\chapter{撒迦利亞書}
\begin{table}[H]
\centering 
\begin{tabular}{p{6.cm}|p{8.5cm}|p{2.5cm}}\hline\hline
波斯王居魯士元年&為建聖殿下詔&拉1:1-4\\\hline
居魯士?年七月初一日起&聚集在耶路撒冷,建築以色列神的壇守住棚節，每日獻祭&拉3:1-7\\\hline
到耶路撒冷第二年二月&興工建造耶和華的殿,匠人立耶和華殿根基&拉3:8-13\\\hline
第二年六月初一日(哈該)&殿荒涼,建造這殿，我就因此喜樂，且得榮耀&哈1:1-12\\\hline
第二年六月二十四日(哈該)&百姓都聽從耶和華的話,耶和華說我與你們同在&哈1:13-15\\\hline
第二年七月二十一日(哈該)&當剛強,我與你們同在;震動天地滄海與旱地萬國,這殿後來的榮耀必大過先前的榮耀&哈2:1-9\\\hline
第二年八月(撒迦利亞)&呼喚百姓不要效法列祖惡行歸向神&\textcolor{red}{撒1:1-6}\\\hline
第二年九月二十四日(哈該)&衣襟兜聖肉;人因摸死屍染了汙穢&哈2:10-19\\\hline
第二年九月二十四日(哈該)&震動天地;以所羅巴伯為印，因我揀選了你&哈2:20-23\\\hline
第二年十一(細罷特)月二十四日&啟示九个異象&撒1:7-6:15\\\hline
第四年九月(基斯流)初四日(撒迦利亞)&責備不誠懇的禁食&撒7:18-14:21\\\hline
第六年十二(亞達)月初三日&猶大長老因先知哈該和易多的孫子撒迦利亞所說勸勉的話，就建造這殿，凡事亨通。他們遵著以色列神的命令和波斯王居魯士、大流士、亞達薛西的旨意，建造完畢。&拉1:14-15\\\hline
\end{tabular}
\caption{波斯王居魯士/大利烏王}
\end{table}

\section{大流士二年,呼喚百姓歸向神,八個異象,為約書亞製金冕示主殿必成 1-6}
\subsection{大利烏王第二年八月,呼喚百姓歸向神 1:1-6 }
大流士王第二年八月，耶和華的話臨到易多的孫子、比利家的兒子先知撒迦利亞說： 2 「耶和華曾向你們列祖大大發怒。 3 所以你要對以色列人說：『萬軍之耶和華如此說：你們要轉向我，我就轉向你們。這是萬軍之耶和華說的。 4 不要效法你們列祖，從前的先知呼叫他們說「萬軍之耶和華如此說：你們要回頭，離開你們的惡道惡行」，他們卻不聽，也不順從我。這是耶和華說的。 5 你們的列祖在哪裡呢？那些先知能永遠存活嗎？ 6 只是我的言語和律例，就是所吩咐我僕人眾先知的，豈不臨到你們列祖嗎？他們就回頭，說：「萬軍之耶和華定意按我們的行動作為向我們怎樣行，他已照樣行了！」』」


\subsection{大流士第二年十一月(細罷特月)二十四日,八個異象 1:7-6:8}
\subsubsection{見乘馬者之異象,錫安再豐盛發達 1:7-17}



\begin{table}[H]
\centering 
\begin{tabular}{p{3.7cm}|p{13.6cm}}\hline\hline
\multicolumn{1}{c|}{撒迦利亞}&\multicolumn{1}{c}{天使}\\\hline
\multirow{3}{3.9cm}{大流士第二年十一月,就是細罷特月二十四日,耶和華的話臨到易多的孫子,比利家的兒子先知撒迦利亞說:我夜間觀看,見一人騎著紅馬,站在窪地番石榴樹中間,在他身後又有紅馬,黃馬和白馬。我對與我說話的天使說:「主啊,這是什麼意思?」}
&他說:「我要指示你這是什麼意思。」那站在番石榴樹中間的人說:「這是奉耶和華差遣在遍地走來走去的」\\\cline{2-2}
&那些騎馬的對站在番石榴樹中間耶和華的使者說:「我們已在遍地走來走去,見全地都安息平靜」於是，耶和華的使者說：「萬軍之耶和華啊，你惱恨耶路撒冷和猶大的城邑已經七十年，你不施憐憫要到幾時呢？」耶和華就用美善的安慰話回答那與我說話的天使。\\\cline{2-2}
&與我說話的天使對我說：「你要宣告說：『萬軍之耶和華如此說：我為耶路撒冷為錫安，心裡極其火熱。我甚惱怒那安逸的列國，因我從前稍微惱怒我民，他們就加害過分。所以耶和華如此說：現今我回到耶路撒冷，仍施憐憫，我的殿必重建在其中，準繩必拉在耶路撒冷之上。這是萬軍之耶和華說的。』你要再宣告說：『萬軍之耶和華如此說：我的城邑必再豐盛發達。耶和華必再安慰錫安，揀選耶路撒冷。』」\\\hline
\end{tabular}
\end{table}

\subsubsection{四角四個匠人的異象 1:18-21}
\begin{table}[H]
\centering 
\begin{tabular}{p{7.cm}|p{8.75cm}}\hline\hline
\multicolumn{1}{c|}{四角}&\multicolumn{1}{c}{四個匠人}\\\hline
我舉目觀看，見有四角。我就問與我說話的天使說：「這是什麼意思？」他回答說：「這是打散猶大、以色列和耶路撒冷的角。」
&耶和華又指四個匠人給我看。我說：「他們來做什麼呢？」他說：「這是打散猶大的角，使人不敢抬頭，但這些匠人來威嚇列國，打掉他們的角，就是舉起打散猶大地的角。」\\\hline
\end{tabular}
\end{table}


\subsubsection{準繩量耶路撒冷,耶和華做火城和其中榮耀的異象 2}
\begin{table}[H]
\centering 
\begin{tabular}{p{2.8cm}|p{14.75cm}}\hline\hline
手拿準繩的人1-3&神做火城和榮耀,錫安民逃脫災害,懲罰列国,再揀選錫安4-13\\\hline
\multirow{4}{3.cm}{我又舉目觀看,見一人手拿準繩。我說:「你往哪裡去?」他對我說:「要去量耶路撒冷,看有多寬多長。」與我說話的天使去的時候,又有一位天使迎著他來，對他說:「你跑去告訴那少年人說:}
&『耶路撒冷必有人居住，好像無城牆的鄉村，因為人民和牲畜甚多。耶和華說：我要做耶路撒冷四圍的火城，並要做其中的榮耀。\\\cline{2-2}
&耶和華說：我從前分散你們在天的四方(猶如天的四風)，現在你們要從北方之地逃回。這是耶和華說的。與巴比倫人同住的錫安民哪，應當逃脫！\\\cline{2-2}
&萬軍之耶和華說,在顯出榮耀之後,差遣我去懲罰那擄掠你們的列國,摸你們的就是摸他眼中的瞳人.看哪我要向他們掄手,他們就必做服侍他們之人的擄物,你們便知道萬軍之耶和華差遣我了\\\cline{2-2}
&『錫安城啊，應當歡樂歌唱！因為我來，要住在你中間。這是耶和華說的。那時，必有許多國歸附耶和華，做他 (我)的子民。他 (我)要住在你中間，你就知道萬軍之耶和華差遣我到你那裡去了。耶和華必收回猶大做他聖地的份，也必再揀選耶路撒冷。凡有血氣的，都當在耶和華面前靜默無聲！因為他興起，從聖所出來了！』」\\\hline
\end{tabular}
%\caption{準繩量耶路撒冷的異象 }
\end{table}


\subsubsection{大祭司約書亞的異象 3}
\begin{table}[H]
\centering
\begin{tabular}{p{11.25cm}|p{6.cm}}\hline\hline
\multicolumn{1}{c|}{耶和華使者}&\multicolumn{1}{c}{約書亞/撒旦/撒迦利亞}\\\hline
天使(他)又指給我看，
&大祭司約書亞站在耶和華的使者面前，撒旦也站在約書亞的右邊，與他作對。\\\hline

耶和華向撒旦說：「撒旦哪，耶和華責備你，就是揀選耶路撒冷的耶和華責備你！這不是從火中抽出來的一根柴嗎？」 
&約書亞穿著汙穢的衣服，站在使者面前。\\\hline
使者吩咐站在面前的說：「你們要脫去他汙穢的衣服。」又對約書亞說：「我使你脫離罪孽，要給你穿上華美的衣服。」 
&我說：「要將潔淨的冠冕戴在他頭上。」\\\hline
他們就把潔淨的冠冕戴在他頭上，給他穿上華美的衣服。耶和華的使者在旁邊站立。耶和華的使者告誡約書亞說： 7 「萬軍之耶和華如此說：你若遵行我的道，謹守我的命令，你就可以管理我的家，看守我的院宇。我也要使你在這些站立的人中間來往。&\multirow{2}{6cm}{}\\\cline{1-1}
大祭司約書亞啊，你和坐在你面前的同伴都當聽（他們是做預兆的）：我必使我僕人大衛的苗裔發出。 9 看哪，我在約書亞面前所立的石頭，在一塊石頭上有七眼。萬軍之耶和華說：我要親自雕刻這石頭，並要在一日之間除掉這地的罪孽。 10 當那日，你們各人要請鄰舍坐在葡萄樹和無花果樹下。這是萬軍之耶和華說的。&\\\hline
\end{tabular}
\end{table}


\subsubsection{金燈台和兩顆橄欖樹的異象,所羅巴伯建殿蒙主佑助 4:1-14}
\begin{table}[H]
\centering
\begin{tabular}{p{8.25cm}|p{9.25cm}}\hline\hline
\multicolumn{1}{c|}{撒迦利亞}&\multicolumn{1}{c}{耶和華使者}\\\hline
那與我說話的天使又來叫醒我好像人睡覺被喚醒一樣
&他問我說：「你看見了什麼？」\\\hline
我說：「我看見了一個純金的燈臺，頂上有燈盞，燈臺上有七盞燈，每盞有七個管子。旁邊有兩棵橄欖樹，一棵在燈盞的右邊，一棵在燈盞的左邊。」我問與我說話的天使說：「主啊，這是什麼意思？」
&與我說話的天使回答我說：「你不知道這是什麼意思嗎？」\\\hline
我說：「主啊，我不知道。」 
&他對我說:「這是耶和華指示所羅巴伯的。萬軍之耶和華說:不是倚靠勢力,不是倚靠才能,乃是倚靠我的靈方能成事,大山哪,你算什麼呢?在所羅巴伯面前,你必成為平地。他必搬出一塊石頭,安在殿頂上,人且大聲歡呼說:『願恩惠恩惠歸於這殿(石)!』」耶和華的話又臨到我說:「所羅巴伯的手立了這殿的根基,他的手也必完成這工,你就知道萬軍之耶和華差遣我到你們這裡來了,誰藐視這日的事為小呢?這七眼乃是耶和華的眼睛,遍察全地,見所羅巴伯手拿線砣就歡喜」\\\hline
我又問天使說：「這燈臺左右的兩棵橄欖樹是什麼意思？」我二次問他說：「這兩根橄欖枝在兩個流出金色油的金嘴旁邊是什麼意思？」&他對我說：「你不知道這是什麼意思嗎？」\\\hline
我說：「主啊，我不知道。」&他說：「這是兩個受膏者，站在普天下主的旁邊。」\\\hline
\end{tabular}
\end{table}

\subsubsection{飛行書卷的異象 5:1-4}
\begin{table}[H]
\centering
\begin{tabular}{p{3.25cm}|p{12.25cm}}\hline\hline
\multicolumn{1}{c|}{撒迦利亞}&\multicolumn{1}{c}{耶和華使者}\\\hline
我又舉目觀看，見有一飛行的書卷。 &他問我說：「你看見了什麼？」\\\hline
我回答說：「我看見一飛行的書卷，長二十肘，寬十肘。」&他對我說：「這是發出行在遍地上的咒詛，凡偷竊的必按卷上這面的話除滅，凡起假誓的必按卷上那面的話除滅。 4 萬軍之耶和華說：我必使這書卷出去，進入偷竊人的家和指我名起假誓人的家，必常在他家裡，連房屋帶木石都毀滅了。」\\\hline
\end{tabular}
\end{table}


\subsubsection{量器中婦人的異象 5:5-11}
\begin{table}[H]
\centering
\begin{tabular}{p{6.75cm}|p{10.cm}}\hline\hline
\multicolumn{1}{c|}{耶和華使者}&\multicolumn{1}{c}{撒迦利亞}\\\hline
與我說話的天使出來，對我說：「你要舉目觀看，見所出來的是什麼。」&我說：「這是什麼呢？」\\\hline
他說：「這出來的是量器。」他又說：「這是惡人在遍地的形狀。」&我見有一片圓鉛被舉起來，這坐在量器中的是個婦人\\\hline
天使說：「這是罪惡。」他就把婦人扔在量器中，將那片圓鉛扔在量器的口上。&我又舉目觀看，見有兩個婦人出來，在她們翅膀中有風，飛得甚快，翅膀如同鸛鳥的翅膀。她們將量器抬起來，懸在天地中間。我問與我說話的天使說：「她們要將量器抬到哪裡去呢？」\\\hline
他對我說：「要往示拿地去,為它蓋造房屋。等房屋齊備,就把它安置在自己的地方。」&\\\hline
\end{tabular}
\end{table}

\subsubsection{四輛馬車出兩座銅山的異象 6:1-8}
\begin{table}[H]
\centering
\begin{tabular}{p{6.75cm}|p{10.cm}}\hline\hline
\multicolumn{1}{c|}{撒迦利亞1-4}&\multicolumn{1}{c}{耶和華使者5-8}\\\hline
我又舉目觀看，見有四輛車從兩山中間出來，那山是銅山。第一輛車套著紅馬，第二輛車套著黑馬，第三輛車套著白馬，第四輛車套著有斑點的壯馬。我就問與我說話的天使說:「主啊，這是什麼意思?」&天使回答我說：「這是天的四風，是從普天下的主面前出來的。套著黑馬的車往北方去，白馬跟隨在後，有斑點的馬往南方去。」壯馬出來，要在遍地走來走去。天使說：「你們只管在遍地走來走去。」牠們就照樣行了。他又呼叫我說：「看哪,往北方去的已在北方安慰我的心。」\\\hline
\end{tabular}
\end{table}
	 
\subsection{約書亞帶金冠冕以示主殿必成的應許 6:9-15}
\begin{table}[H]
\centering
\begin{tabular}{p{5.45cm}|p{5.35cm}|p{6.2cm}}\hline\hline
\multicolumn{1}{c|}{取金銀做冠冕9-11}&\multicolumn{1}{c|}{大衛苗裔12-13}&\multicolumn{1}{c}{冠冕歸約西亞放在殿裡為紀念14-15}\\\hline
耶和華的話臨到我說:「你要從被擄之人中取黑玳、多比雅、耶大雅的金銀,這三人是從巴比倫來到西番雅的兒子約西亞的家裡。當日你要進他的家,取這金銀做冠冕,戴在約薩答的兒子大祭司約書亞的頭上
&對他說:『萬軍之耶和華如此說:看哪,那名稱為大衛苗裔的,他要在本處長起來,並要建造耶和華的殿。他要建造耶和華的殿,並擔負尊榮,坐在位上掌王權,又必在位上做祭司,使兩職之間籌定和平。』
&這冠冕要歸希連(黑玳),多比雅,耶大雅和西番雅的兒子賢(約西亞),放在耶和華的殿裡為紀念。遠方的人也要來建造耶和華的殿,你們就知道萬軍之耶和華差遣我到你們這裡來。你們若留意聽從耶和華你們神的話,這事必然成就。」\\\hline
\end{tabular}
\end{table}



\section{神的應許 7-14}
\subsection{禁食的日子必變為猶大家歡喜快樂的日子和節期 7-8}

\subsubsection{被擄者詢主仍禁食否 7:1-14}
\begin{table}[H]
\centering
\begin{tabular}{p{4.25cm}|p{4.55cm}|p{8.2cm}}\hline\hline
\multicolumn{1}{c|}{被擄者詢主仍禁食否1-3}&\multicolumn{1}{c|}{責備不敬虔的禁食4-7}&\multicolumn{1}{c}{众民因罪被擄8-14}\\\hline
大流士王第四年九月，就是基斯流月初四日，耶和華的話臨到撒迦利亞。那時伯特利人已經打發沙利色和利堅米勒並跟從他們的人，去懇求耶和華的恩，並問萬軍之耶和華殿中的祭司和先知說：「我歷年以來，在五月間哭泣齋戒，現在還當這樣行嗎？」
&萬軍之耶和華的話就臨到我說:「你要宣告國內的眾民和祭司說:你們這七十年在五月、七月禁食悲哀,豈是絲毫向我禁食嗎?你們吃喝,不是為自己吃、為自己喝嗎?當耶路撒冷和四圍的城邑有居民，正興盛,南地高原有人居住的時候,耶和華藉從前的先知所宣告的話，你們不當聽嗎?」
&耶和華的話又臨到撒迦利亞說：「萬軍之耶和華曾對你們的列祖如此說：『要按至理判斷，各人以慈愛憐憫弟兄。不可欺壓寡婦、孤兒、寄居的和貧窮人，誰都不可心裡謀害弟兄。』他們卻不肯聽從，扭轉肩頭，塞耳不聽，使心硬如金鋼石，不聽律法和萬軍之耶和華用靈藉從前的先知所說的話。故此，萬軍之耶和華大發烈怒。萬軍之耶和華說：我曾呼喚他們，他們不聽，將來他們呼求我，我也不聽。我必以旋風吹散他們到素不認識的萬國中。這樣，他們的地就荒涼，甚至無人來往經過，因為他們使美好之地荒涼了。」\\\hline
\end{tabular}
\end{table}


\subsubsection{耶路撒冷必復振興 8:1-8}
%\subsubsubsection{耶路撒冷必復振興 8:1-8} 
\begin{itemize} 
\itemsep-.5em 
\item 萬軍之耶和華的話臨到我說：「萬軍之耶和華如此說：我為錫安心裡極其火熱，我為她火熱，向她的仇敵發烈怒。
\item 耶和華如此說：我現在回到錫安，要住在耶路撒冷中。耶路撒冷必稱為誠實的城，萬軍之耶和華的山必稱為聖山。
\item 萬軍之耶和華如此說：將來必有年老的男女坐在耶路撒冷街上，因為年紀老邁就手拿拐杖。城中街上必滿有男孩女孩玩耍。
\item 萬軍之耶和華如此說：到那日，這事在餘剩的民眼中看為稀奇，在我眼中也看為稀奇嗎？這是萬軍之耶和華說的。
\item 萬軍之耶和華如此說：我要從東方從西方救回我的民。我要領他們來，使他們住在耶路撒冷中。他們要做我的子民，我要做他們的神，都憑誠實和公義。
\end{itemize}

\subsubsection{勖建殿者不要懼怕 8:9-13} 
\begin{itemize} 
\itemsep-.5em 
\item 萬軍之耶和華如此說：當建造萬軍之耶和華的殿立根基之日的先知所說的話，現在你們聽見，應當手裡強壯。 
\item 那日以先，人得不著雇價，牲畜也是如此；且因敵人的緣故，出入之人不得平安，乃因我使眾人互相攻擊。 
\item 但如今，我待這餘剩的民必不像從前。這是萬軍之耶和華說的。因為他們必平安撒種，葡萄樹必結果子，地土必有出產，天也必降甘露。我要使這餘剩的民享受這一切的福。
\item 猶大家和以色列家啊，你們從前在列國中怎樣成為可咒詛的，照樣，我要拯救你們，使人稱你們為有福的 (使你們叫人得福)。你們不要懼怕，手要強壯！
\end{itemize}
萬軍之耶和華如此說：當建造萬軍之耶和華的殿立根基之日的先知所說的話，現在你們聽見，應當手裡強壯。 10 那日以先，人得不著雇價，牲畜也是如此；且因敵人的緣故，出入之人不得平安，乃因我使眾人互相攻擊。 11 但如今，我待這餘剩的民必不像從前。這是萬軍之耶和華說的。 12 因為他們必平安撒種，葡萄樹必結果子，地土必有出產，天也必降甘露。我要使這餘剩的民享受這一切的福。 13 猶大家和以色列家啊，你們從前在列國中怎樣成為可咒詛的，照樣，我要拯救你們，使人稱你們為有福的（使你們叫人得福）。你們不要懼怕，手要強壯！

\subsubsection{勉以善行 8:14-17} 
\begin{itemize} 
\itemsep-.5em 
\item 萬軍之耶和華如此說：你們列祖惹我發怒的時候，我怎樣定意降禍並不後悔， 現在我照樣定意施恩於耶路撒冷和猶大家
\item 你們不要懼怕。 你們所當行的是這樣：各人與鄰舍說話誠實，在城門口按至理判斷，使人和睦。誰都不可心裡謀害鄰舍，也不可喜愛起假誓，因為這些事都為我所恨惡。這是耶和華說的。
\end{itemize}
萬軍之耶和華如此說：你們列祖惹我發怒的時候，我怎樣定意降禍並不後悔， 15 現在我照樣定意施恩於耶路撒冷和猶大家，你們不要懼怕。 16 你們所當行的是這樣：各人與鄰舍說話誠實，在城門口按至理判斷，使人和睦。 17 誰都不可心裡謀害鄰舍，也不可喜愛起假誓，因為這些事都為我所恨惡。這是耶和華說的。

\subsubsection{憂傷將易歡欣 8:18-19}
萬軍之耶和華的話臨到我說： 19 「萬軍之耶和華如此說：四月、五月禁食的日子，七月、十月禁食的日子，必變為猶大家歡喜快樂的日子和歡樂的節期，所以你們要喜愛誠實與和平。

\subsubsection{萬民必尋求主恩 8:20-23}
\begin{itemize} 
\itemsep-.5em 
\item 萬軍之耶和華如此說：將來必有列國的人和多城的居民來到。這城的居民必到那城說：『我們要快去懇求耶和華的恩，尋求萬軍之耶和華！我也要去！』必有列邦的人和強國的民來到耶路撒冷，尋求萬軍之耶和華，懇求耶和華的恩。
\item 萬軍之耶和華如此說：在那些日子，必有十個人從列國諸族(方言)中出來，拉住一個猶大人的衣襟，說：『我們要與你們同去，因為我們聽見神與你們同在了。
\end{itemize}
萬軍之耶和華如此說：將來必有列國的人和多城的居民來到。 21 這城的居民必到那城說：『我們要快去懇求耶和華的恩，尋求萬軍之耶和華！我也要去！』 22 必有列邦的人和強國的民來到耶路撒冷，尋求萬軍之耶和華，懇求耶和華的恩。 23 萬軍之耶和華如此說：在那些日子，必有十個人從列國諸族 （方言）中出來，拉住一個猶大人的衣襟，說：『我們要與你們同去，因為我們聽見神與你們同在了。』



\subsection{審判列國佑其民使之獲勝的應許 9}
\subsubsection{審判列國 9:1-7}
\begin{table}[!hbp]
\centering 
\begin{tabular}{|p{3.3cm}||p{12cm}|}
\hline
\hline
哈得拉地大馬色 &耶和華的默示應驗在哈得拉地大馬色。世人和以色列各支派的眼目都仰望耶和華\\\hline
靠近的哈馬 &耶和華的默示應驗在靠近的哈馬。世人和以色列各支派的眼目都仰望耶和華\\\hline
推羅、西頓 &耶和華的默示應驗在推羅、西頓。世人和以色列各支派的眼目都仰望耶和華。因為這二城的人大有智慧。推羅為自己修築保障，積蓄銀子如塵沙，堆起精金如街上的泥土。主必趕出他，打敗他海上的權利；他必被火燒滅。\\\hline
亞實基倫& 看見必懼怕，也不再有居民。\\\hline
迦薩&看見甚痛苦,必不再有君王；\\\hline
 以革倫 &因失了盼望蒙羞；以革倫人必如耶布斯人\\\hline
 亞實突&私生子（外族人）必住在亞實突\\\hline
 非利士人&我必除滅非利士人的驕傲。我必除去他口中帶血之肉和牙齒內可憎之物。他必作為餘剩的人歸與我們的神，必在猶大像族長；\\\hline
\end{tabular}
\caption{審判列國 }
\end{table}
耶和華的默示應驗在哈得拉地大馬士革——世人和以色列各支派的眼目都仰望耶和華—— 2 和靠近的哈馬，並推羅、西頓，因為這二城的人大有智慧。 3 「推羅為自己修築保障，積蓄銀子如塵沙，堆起精金如街上的泥土。 4 主必趕出她，打敗她海上的權力，她必被火燒滅。 5 亞實基倫看見必懼怕，加沙看見甚痛苦，以革倫因失了盼望蒙羞。加沙必不再有君王，亞實基倫也不再有居民。 6 私生子（外族人）必住在亞實突，我必除滅非利士人的驕傲。 7 我必除去他口中帶血之肉和牙齒內可憎之物，他必作為餘剩的人歸於我們的神，必在猶大像族長，以革倫人必如耶布斯人。




\subsubsection{看顧錫安的應許 9:8-17}
\begin{itemize} 
\itemsep-.5em 
\item 我必在我家的四圍安營，使敵軍不得任意往來，暴虐的人也不再經過，因為我親眼看顧我的家
\item 和平再临錫安
\begin{itemize} 
\itemsep-.5em 
\item 錫安的民哪，應當大大喜樂！耶路撒冷的民哪，應當歡呼！看哪，你的王來到你這裡，他是公義的，並且施行拯救，謙謙和和地騎著驢，就是騎著驢的駒子。
\item 我必除滅以法蓮的戰車和耶路撒冷的戰馬，爭戰的弓也必除滅。他必向列國講和平，他的權柄必從這海管到那海，從大河管到地極
\end{itemize}
\item 主佑其民使之獲勝
\begin{itemize} 
\itemsep-.5em 
\item 錫安哪，我因與你立約的血，將你中間被擄而囚的人從無水的坑中釋放出來。 12 你們被囚而有指望的人，都要轉回保障！我今日說明：我必加倍賜福給你們。
\item 我拿猶大做上弦的弓，我拿以法蓮為張弓的箭。錫安哪，我要激發你的眾子攻擊希臘（雅完）的眾子，使你如勇士的刀。
\item 耶和華必顯現在他們以上，他的箭必射出像閃電。主耶和華必吹角，乘南方的旋風而行。 萬軍之耶和華必保護他們。他們必吞滅仇敵，踐踏彈石；他們必喝血呐喊，猶如飲酒；他們必像盛滿血的碗，又像壇的四角滿了血。
\item 當那日，耶和華他們的神必看他的民如群羊，拯救他們。因為他們必像冠冕上的寶石，高舉在他的地以上（在他的地上發光輝）。 
\item 他的恩慈何等大，他的榮美何其盛！五穀健壯少男，新酒培養處女
\end{itemize}
\end{itemize}
8 「我必在我家的四圍安營，使敵軍不得任意往來，暴虐的人也不再經過，因為我親眼看顧我的家。錫安的民哪，應當大大喜樂！耶路撒冷的民哪，應當歡呼！看哪，你的王來到你這裡，他是公義的，並且施行拯救，謙謙和和地騎著驢，就是騎著驢的駒子。 10 我必除滅以法蓮的戰車和耶路撒冷的戰馬，爭戰的弓也必除滅。他必向列國講和平，他的權柄必從這海管到那海，從大河管到地極。
錫安哪，我因與你立約的血，將你中間被擄而囚的人從無水的坑中釋放出來。 12 你們被囚而有指望的人，都要轉回保障！我今日說明：我必加倍賜福給你們。 13 我拿猶大做上弦的弓，我拿以法蓮為張弓的箭。錫安哪，我要激發你的眾子攻擊希臘（雅完）的眾子，使你如勇士的刀。 14 耶和華必顯現在他們以上，他的箭必射出像閃電。主耶和華必吹角，乘南方的旋風而行。 15 萬軍之耶和華必保護他們。他們必吞滅仇敵，踐踏彈石；他們必喝血呐喊，猶如飲酒；他們必像盛滿血的碗，又像壇的四角滿了血。 16 當\textcolor{red}{那日}，耶和華他們的神必看他的民如群羊，拯救他們。因為他們必像冠冕上的寶石，高舉在他的地以上(在他的地上發光輝)。 17 他的恩慈何等大，他的榮美何其盛！五穀健壯少男，新酒培養處女。

\subsection{责备假先知牧人，与百姓同在重返故土的應許 10-11}
\subsubsection{责备假先知牧人 10:1-5}
「當春雨的時候，你們要向發閃電的耶和華求雨，他必為眾人降下甘霖，使田園生長菜蔬。 2 因為家神所言的是虛空，卜士所見的是虛假，做夢者所說的是假夢。他們白白地安慰人，所以眾人如羊流離，因無牧人就受苦。 

3 我的怒氣向牧人發作，我必懲罰公山羊，因我萬軍之耶和華眷顧自己的羊群，就是猶大家，必使他們如駿馬在陣上。 4 房角石、釘子、爭戰的弓和一切掌權的都從他而出。 5 他們必如勇士，在陣上將仇敵踐踏在街上的泥土中。他們必爭戰，因為耶和華與他們同在，騎馬的也必羞愧。


\subsubsection{主必领民使返故土的應許 10:6-12}

我要堅固猶大家，拯救約瑟家，要領他們歸回。我要憐恤他們，他們必像未曾棄絕的一樣。都因我是耶和華他們的神，我必應允他們的禱告。 7 以法蓮人必如勇士，他們心中暢快如同喝酒。他們的兒女必看見而快活，他們的心必因耶和華喜樂。

8 「我要發嘶聲，聚集他們，因我已經救贖他們。他們的人數必加增，如從前加增一樣。 9 我雖然（必）播散他們在列國中，他們必在遠方記念我。他們與兒女都必存活，且得歸回。 10 我必再領他們出埃及地，招聚他們出亞述，領他們到基列和黎巴嫩，這地尚且不夠他們居住。 11 耶和華必經過苦海，擊打海浪，使尼羅河的深處都枯乾。亞述的驕傲必致卑微，埃及的權柄必然滅沒。 12 我必使他們倚靠我得以堅固，一舉一動必奉我的名。這是耶和華說的。


\subsubsection{牧人哀號草场荒廢 11:1-3}
黎巴嫩哪，開開你的門，任火燒滅你的香柏樹！ 2 松樹啊，應當哀號！因為香柏樹傾倒，佳美的樹毀壞。巴珊的橡樹啊，應當哀號！因為茂盛的樹林已經倒了。 3 聽啊！有牧人哀號的聲音，因他們榮華的草場毀壞了；有少壯獅子咆哮的聲音，因約旦河旁的叢林荒廢了


\subsubsection{以二杖牧群羊 11:4-9}
耶和華我的神如此說：「你撒迦利亞要牧養這將宰的群羊。 5 買牠們的宰了牠們，以自己為無罪。賣牠們的說：『耶和華是應當稱頌的！因我成為富足。』牧養牠們的並不憐恤牠們。 6 耶和華說：我不再憐恤這地的居民，必將這民交給各人的鄰舍和他們王的手中。他們必毀滅這地，我也不救這民脫離他們的手。」 7 於是，我牧養這將宰的群羊，就是群中最困苦的羊。我拿著兩根杖，一根我稱為「榮美」，一根我稱為「聯索」。這樣，我牧養了群羊。 8 「一月之內，我除滅三個牧人，因為我的心厭煩他們，他們的心也憎嫌我。 9 我就說：『我不牧養你們。要死的由他死，要喪亡的由他喪亡，餘剩的由他們彼此相食。』

%\subsection{折斷二杖表明廢約 11:10-17}
\subsubsection{折斷二杖 11:10-14}
我折斷那稱為「榮美」的杖，表明我廢棄與萬民所立的約。 11 當日就廢棄了。這樣，那些仰望我的困苦羊，就知道所說的是耶和華的話。 12 我對他們說：「你們若以為美，就給我工價；不然，就罷了。」於是他們給了三十塊錢，作為我的工價。 13 耶和華吩咐我說：「要把眾人所估定美好的價值丟給窯戶。」我便將這三十塊錢，在耶和華的殿中丟給窯戶了。
14 我又折斷稱為「聯索」的那根杖，表明我廢棄猶大與以色列弟兄的情誼。

\subsubsection{興起一個坏牧人 11:15-17}
15 耶和華又吩咐我說：「你再取愚昧人所用的器具。 16 因我要在這地興起一個牧人，他不看顧喪亡的，不尋找分散的，不醫治受傷的，也不牧養強壯的，卻要吃肥羊的肉，撕裂牠的蹄子。 17 無用的牧人丟棄羊群有禍了！刀必臨到他的膀臂和右眼上，他的膀臂必全然枯乾，他的右眼也必昏暗失明。

\subsection{洗除耶路撒冷罪惡與汙穢，以色列，猶大復興的默示 12-13}
\subsubsection{以耶路撒冷為杯使四周之民致醉昏眩 12:1-4}
耶和華論以色列的默示。鋪張諸天，建立地基，造人裡面之靈的耶和華說： 2 「我必使耶路撒冷被圍困的時候，向四圍列國的民成為令人昏醉的杯，這默示也論到猶大 （猶大也是如此）。\textcolor{red}{那日}，我必使耶路撒冷向聚集攻擊她的萬民當做一塊重石頭，凡舉起的必受重傷。」 4 耶和華說：「\textcolor{red}{到那日}，我必使一切馬匹驚惶，使騎馬的顛狂。我必看顧猶大家，使列國的一切馬匹瞎眼。

\subsubsection{許猶大復興,滅絕來攻擊耶路撒冷各國的民 12:5-9}
 猶大的族長必心裡說：『耶路撒冷的居民倚靠萬軍之耶和華他們的神，就做我們的能力。』 6 \textcolor{red}{那日}，我必使猶大的族長如火盆在木柴中，又如火把在禾捆裡，他們必左右燒滅四圍列國的民。耶路撒冷人必仍住本處，就是耶路撒冷。 7 耶和華必先拯救猶大的帳篷，免得大衛家的榮耀和耶路撒冷居民的榮耀勝過猶大。 8 \textcolor{red}{那日}，耶和華必保護耶路撒冷的居民，他們中間軟弱的必如大衛，大衛的家必如神，如行在他們前面之耶和華的使者。 9 \textcolor{red}{那日}，我必定意滅絕來攻擊耶路撒冷各國的民。
\subsubsection{耶路撒冷居民必為其所刺者哀哭 12:10-14}
「我必將那施恩叫人懇求的靈澆灌大衛家和耶路撒冷的居民。他們必仰望我(他)，就是他們所扎的；必為我悲哀，如喪獨生子；又為我愁苦，如喪長子。 11 \textcolor{red}{那日}，耶路撒冷必有大大的悲哀，如米吉多平原之哈達臨門的悲哀。 12 境內一家一家的都必悲哀。\textcolor{blue}{大衛}家，男的獨在一處，女的獨在一處。\textcolor{blue}{拿單}家，男的獨在一處，女的獨在一處。 13 \textcolor{blue}{利未}家，男的獨在一處，女的獨在一處。\textcolor{blue}{示每}家，男的獨在一處，女的獨在一處。 14 \textcolor{blue}{其餘的}各家，男的獨在一處，女的獨在一處。


%\subsection{洗除耶路撒冷罪惡與汙穢 13}
\subsubsection{闢一源洗罪泉源,滅絕偶像偽先知 13:1-6}
\begin{table}[H]
\centering
\begin{tabular}{p{2.25cm}|p{14.75cm}}\hline\hline
\multicolumn{1}{c|}{洗罪泉源1}&\multicolumn{1}{c}{滅絕偶像偽先知2-6}\\\hline
\textcolor{red}{那日}，必給大衛家和耶路撒冷的居民開一個泉源，洗除罪惡與汙穢。
&萬軍之耶和華說：「\textcolor{red}{那日}，我必從地上除滅偶像的名，不再被人記念，也必使這地不再有假先知與汙穢的靈。若再有人說預言，生他的父母必對他說：『你不得存活，因為你託耶和華的名說假預言。』生他的父母在他說預言的時候，要將他刺透。\textcolor{red}{那日}，凡做先知說預言的必因他所論的異象羞愧，不再穿毛衣哄騙人。他必說：『我不是先知，我是耕地的，我從幼年做人的奴僕。』必有人問他說：『你兩臂中間是什麼傷呢？』他必回答說：『這是我在親友家中所受的傷。』」\\\hline
\end{tabular}
\end{table}


\subsubsection{牧人被擊,全地三分之一归神 13:7-9}

萬軍之耶和華說：「刀劍哪，應當興起，攻擊我的牧人和我的同伴！擊打牧人，羊就分散，我必反手加在微小者的身上。」 

8 耶和華說：「這全地的人，三分之二必剪除而死，三分之一仍必存留。 9 我要使這三分之一經火，熬煉他們如熬煉銀子，試煉他們如試煉金子。他們必求告我的名，我必應允他們。我要說：『這是我的子民。』他們也要說：『耶和華是我們的神。』」

\subsection{主足必臨於橄欖山,耶和華為全地之王 14}
\subsubsection{萬國攻打耶路撒冷,城中一半被擄，剩下的在城中不致剪除 14:1-3}
耶和華的日子臨近，你的\textcolor{red}{財物}必被搶掠，在你中間分散。 2 「因為我必聚集萬國與耶路撒冷爭戰，\textcolor{red}{城}必被攻取，\textcolor{red}{房屋}被搶奪，\textcolor{red}{婦女}被玷汙，城中的\textcolor{red}{民一半}被擄去，\textcolor{red}{剩下的民}仍在城中，不致剪除。」 3 那時耶和華必出去與那些國爭戰，好像從前爭戰一樣。

\subsubsection{主足必臨於橄欖山 14:4-11}
\begin{table}[H]
\centering
\begin{tabular}{p{6.25cm}|p{2.250cm}|p{2.cm}|p{6.250cm}}\hline\hline
\multicolumn{1}{c|}{主足必臨於橄欖山4-5}&\multicolumn{1}{c|}{三光退縮6-7}&\multicolumn{1}{c|}{活水8}&\multicolumn{1}{c}{耶和華為全地之王9-11}\\\hline
\textcolor{red}{那日},他的腳必站在耶路撒冷前面朝東的橄欖山上。這山必從中間分裂,自東至西成為極大的谷,山的一半向北挪移,一半向南挪移。「你們要從我山的谷中逃跑,因為山谷必延到亞薩。你們逃跑,必如猶大王烏西雅年間的人逃避大地震一樣。」耶和華我的神必降臨,有一切聖者同來。
&\textcolor{red}{那日}必沒有光，三光必退縮。\textcolor{red}{那日}必是耶和華所知道的，不是白晝，也不是黑夜，到了晚上才有光明。
&\textcolor{red}{那日}必有活水從耶路撒冷出來，一半往東海流，一半往西海流，冬夏都是如此。
&耶和華必做全地的王,\textcolor{red}{那日}耶和華必為獨一無二的,他的名也是獨一無二的。全地,從迦巴直到耶路撒冷南方的臨門要變為亞拉巴。耶路撒冷必仍居高位,就是從便雅憫門到第一門之處,又到角門,並從哈楠業樓直到王的酒榨。人必住在其中,不再有咒詛,耶路撒冷人必安然居住。
\\\hline
\end{tabular}
\end{table}

\subsubsection{攻擊耶路撒冷列國人的災殃,其中剩下的人年年來耶路撒冷 14:12-21}
\begin{table}[H]
\centering
\begin{tabular}{p{6.5cm}|p{11cm}}\hline\hline
\multicolumn{1}{c|}{攻擊耶路撒冷列國人的災殃12-15}&\multicolumn{1}{c}{剩下的人年年來耶路撒冷16-21}\\\hline
\multirow{2}{6.5cm}{耶和華用災殃攻擊那與耶路撒冷爭戰的列國人，必是這樣：他們兩腳站立的時候，肉必消沒，眼在眶中乾癟，舌在口中潰爛。\textcolor{red}{那日}，耶和華必使他們大大擾亂，他們各人彼此揪住，舉手攻擊。猶大也必在耶路撒冷爭戰。那時四圍各國的財物，就是許多金銀、衣服，必被收聚。那臨到馬匹、騾子、駱駝、驢和營中一切牲畜的災殃是與那災殃一般。}
&所有來攻擊耶路撒冷列國中剩下的人，必年年上來敬拜大君王萬軍之耶和華，並守住棚節。地上萬族中，凡不上耶路撒冷敬拜大君王萬軍之耶和華的，必無雨降在他們的地上。埃及族若不上來，雨也不降在他們的地上。凡不上來守住棚節的列國人，耶和華也必用這災攻擊他們。這就是埃及的刑罰和那不上來守住棚節之列國的刑罰。\\\cline{2-2}
&當\textcolor{red}{那日}馬的鈴鐺上必有「歸耶和華為聖的」這句話。耶和華殿內的鍋必如祭壇前的碗一樣。凡耶路撒冷和猶大的鍋都必歸萬軍之耶和華為聖,凡獻祭的都必來取這鍋,煮肉在其中。當\textcolor{red}{那日},在萬軍之耶和華的殿中必不再有迦南人。\\\hline
\end{tabular}
\end{table}






